% GENERATED FILE. Edit canonical lessons, then run npm run generate:books.
% canonical-source-hash: fnv1a64:14440ce770ce8f98
% canonical-lessons: MR-C199-cikatavne, MR-C199-gathne, MR-C199-bhet-dene, MR-C199-vinne, MR-C199-kapni-karne

\chapter{To stick, To reach, To visit, To knit, To harvest}
\label{ch:mr-to-stick-to-reach-to-visit-to-knit-to-harvest}
\hlchaptermodality{199}

\begin{chapteropening}
\textbf{By the end of this chapter:} \emph{I can say to stick, to reach, to visit, to knit and to harvest.}

Complete the last lesson of chapter 199: I can say to stick, to reach, to visit, to knit and to harvest.
\end{chapteropening}

\section[\texorpdfstring{\mr{चिकटवणे} (cikaṭavṇe)}{चिकटवणे (cikaṭavṇe)}]{\mr{चिकटवणे} (cikaṭavṇe) — to stick, to glue}

\label{lesson:MR-C199-cikatavne}



\begin{quote}
Before the new one: say the Marathi for to vote, then the Marathi for to welcome.
\end{quote}

\subsection*{You'll want to know: \mr{चिकटवणे}}
\textbf{\mr{चिकटवणे}} — \emph{cikaṭavṇe} — \textquotedblleft{}to stick, to glue\textquotedblright{}.

\begin{grammarlens}[title={What you've built}]
One more word for this chapter.
\end{grammarlens}

\subsection*{Your turn}
\begin{itemize}
\raggedright
  \item \emph{Say it:} \emph{cikaṭavṇe}
  \item \emph{Say it:} \emph{cikaṭavṇe}, once more
  \item \emph{Say it:} say \emph{cikaṭavṇe}
  \item \emph{Recall:} say the Marathi for to vote, then the Marathi for to welcome, then say \emph{cikaṭavṇe} again
\end{itemize}

\begin{tcolorbox}[breakable,colback=teal!4,colframe=teal!35!black,title={Before you move on}]
What does \mr{चिकटवणे} mean? (\textquotedblleft{}To stick, to glue\textquotedblright{}.) Say it once more.
\end{tcolorbox}
\section[\texorpdfstring{\mr{गाठणे} (gāṭhṇe)}{गाठणे (gāṭhṇe)}]{\mr{गाठणे} (gāṭhṇe) — to reach, to catch up with}

\label{lesson:MR-C199-gathne}



\begin{quote}
Before the new one: say the Marathi for to welcome, then the Marathi for to stick.
\end{quote}

\subsection*{You'll want to know: \mr{गाठणे}}
\textbf{\mr{गाठणे}} — \emph{gāṭhṇe} — \textquotedblleft{}to reach, to catch up with\textquotedblright{}.

\begin{grammarlens}[title={What you've built}]
One more word for this chapter.
\end{grammarlens}

\subsection*{Your turn}
\begin{itemize}
\raggedright
  \item \emph{Say it:} \emph{gāṭhṇe}
  \item \emph{Say it:} \emph{gāṭhṇe}, once more
  \item \emph{Say it:} say \emph{gāṭhṇe}
  \item \emph{Recall:} say the Marathi for to welcome, then the Marathi for to stick, then say \emph{gāṭhṇe} again
\end{itemize}

\begin{tcolorbox}[breakable,colback=teal!4,colframe=teal!35!black,title={Before you move on}]
What does \mr{गाठणे} mean? (\textquotedblleft{}To reach, to catch up with\textquotedblright{}.) Say it once more.
\end{tcolorbox}
\section[\texorpdfstring{\mr{भेट} \mr{देणे} (bheṭ deṇe)}{भेट देणे (bheṭ deṇe)}]{\mr{भेट} \mr{देणे} (bheṭ deṇe) — to visit}

\label{lesson:MR-C199-bhet-dene}



\begin{quote}
Before the new one: say the Marathi for to stick, then the Marathi for to reach.
\end{quote}

\subsection*{You'll want to know: \mr{भेट} \mr{देणे}}
\textbf{\mr{भेट} \mr{देणे}} — \emph{bheṭ deṇe} — \textquotedblleft{}to visit\textquotedblright{}.

\begin{grammarlens}[title={What you've built}]
One more word for this chapter.
\end{grammarlens}

\subsection*{Your turn}
\begin{itemize}
\raggedright
  \item \emph{Say it:} \emph{bheṭ deṇe}
  \item \emph{Say it:} \emph{bheṭ deṇe}, once more
  \item \emph{Say it:} say \emph{bheṭ deṇe}
  \item \emph{Recall:} say the Marathi for to stick, then the Marathi for to reach, then say \emph{bheṭ deṇe} again
\end{itemize}

\begin{tcolorbox}[breakable,colback=teal!4,colframe=teal!35!black,title={Before you move on}]
What does \mr{भेट} \mr{देणे} mean? (\textquotedblleft{}To visit\textquotedblright{}.) Say it once more.
\end{tcolorbox}
\section[\texorpdfstring{\mr{विणणे} (viṇṇe)}{विणणे (viṇṇe)}]{\mr{विणणे} (viṇṇe) — to knit, to weave}

\label{lesson:MR-C199-vinne}



\begin{quote}
Before the new one: say the Marathi for to reach, then the Marathi for to visit.
\end{quote}

\subsection*{You'll want to know: \mr{विणणे}}
\textbf{\mr{विणणे}} — \emph{viṇṇe} — \textquotedblleft{}to knit, to weave\textquotedblright{}.

\begin{grammarlens}[title={What you've built}]
One more word for this chapter.
\end{grammarlens}

\subsection*{Your turn}
\begin{itemize}
\raggedright
  \item \emph{Say it:} \emph{viṇṇe}
  \item \emph{Say it:} \emph{viṇṇe}, once more
  \item \emph{Say it:} say \emph{viṇṇe}
  \item \emph{Recall:} say the Marathi for to reach, then the Marathi for to visit, then say \emph{viṇṇe} again
\end{itemize}

\begin{tcolorbox}[breakable,colback=teal!4,colframe=teal!35!black,title={Before you move on}]
What does \mr{विणणे} mean? (\textquotedblleft{}To knit, to weave\textquotedblright{}.) Say it once more.
\end{tcolorbox}
\section[\texorpdfstring{\mr{कापणी} \mr{करणे} (kāpṇī karṇe)}{कापणी करणे (kāpṇī karṇe)}]{\mr{कापणी} \mr{करणे} (kāpṇī karṇe) — to harvest}

\label{lesson:MR-C199-kapni-karne}



\begin{quote}
Before the new one: say the Marathi for to visit, then the Marathi for to knit.
\end{quote}

\subsection*{You'll want to know: \mr{कापणी} \mr{करणे}}
\textbf{\mr{कापणी} \mr{करणे}} — \emph{kāpṇī karṇe} — \textquotedblleft{}to harvest\textquotedblright{}.

\begin{grammarlens}[title={What you've built}]
One more word for this chapter.
\end{grammarlens}

\subsection*{Your turn}
\begin{itemize}
\raggedright
  \item \emph{Say it:} \emph{kāpṇī karṇe}
  \item \emph{Say it:} \emph{kāpṇī karṇe}, once more
  \item \emph{Say it:} say \emph{kāpṇī karṇe}
  \item \emph{Recall:} say the Marathi for to visit, then the Marathi for to knit, then say \emph{kāpṇī karṇe} again
\end{itemize}

\begin{tcolorbox}[breakable,colback=teal!4,colframe=teal!35!black,title={Before you move on}]
What does \mr{कापणी} \mr{करणे} mean? (\textquotedblleft{}To harvest\textquotedblright{}.) Say it once more.
\end{tcolorbox}
